% Generated draft: requires mathematical and source review.
\documentclass[11pt]{article}
\usepackage[margin=25mm]{geometry}
\usepackage{fontspec}
\setmainfont{Latin Modern Roman}
\usepackage{amsmath,amssymb,mathtools,unicode-math}
\setmathfont{Latin Modern Math}
\usepackage{enumitem,needspace,xurl,hyperref}
\setlength{\parindent}{0pt}
\setlength{\parskip}{5pt}
\setlength{\emergencystretch}{3em}
\newcommand{\N}{\mathbb{N}}
\newcommand{\Z}{\mathbb{Z}}
\newcommand{\Q}{\mathbb{Q}}
\newcommand{\R}{\mathbb{R}}
\newcommand{\C}{\mathbb{C}}
\newcommand{\F}{\mathbb{F}}
\newcommand{\A}{\mathbb{A}}
\newcommand{\PP}{\mathbb{P}}
\newcommand{\E}{\mathbb{E}}
\newcommand{\eps}{\varepsilon}
\newcommand{\dd}{\,\mathrm{d}}
\begin{document}
\section*{UnsolvedMath ID: 801158}
\textbf{Source problem code:} catalogue entry
\subsection{Definitions and mathematical statement}
For orientation, a conventional core definition of a Frobenius manifold consists of a manifold \(M\), a nondegenerate symmetric metric \(g\), a commutative, associative multiplication \(\circ\) on the tangent bundle, and a unit vector field \(e\), satisfying
\[
e\circ X=X,\qquad g(X\circ Y,Z)=g(X,Y\circ Z)
\]
for all vector fields \(X,Y,Z\). If \(\nabla\) is the Levi--Civita connection of \(g\), one requires \(g\) to be flat, \(\nabla e=0\), and the covariant tensor \(\nabla c\) to be totally symmetric, where
\[
c(X,Y,Z)=g(X\circ Y,Z).
\]
This is explanatory terminology rather than a determination of the source's complete axiom system: the supplied record does not specify the real or complex category, metric signature, dimension, or whether Euler-field and homogeneity axioms are included.

\textbf{Source problem.} Determine whether Frobenius structure is projectively invariant, and whether its defining flat metric can be replaced by a metric of constant curvature.

For reference, two torsion-free connections \(\nabla\) and \(\widehat\nabla\) are projectively equivalent when they have the same unparametrized geodesics; locally this is expressed by
\[
\widehat\nabla_XY=\nabla_XY+\alpha(X)Y+\alpha(Y)X
\]
for a one-form \(\alpha\). The supplied record does not specify whether this equivalence of connections is the intended meaning of projective invariance, which projectively equivalent metrics are allowed, or how \(\circ\), \(e\), and any additional Frobenius data should transform. Consequently, it does not determine a precise universal invariance assertion.

With the convention
\[
R^g(X,Y)Z=\nabla_X\nabla_YZ-\nabla_Y\nabla_XZ-\nabla_{[X,Y]}Z,
\]
the usual constant-curvature replacement is
\[
R^g(X,Y)Z=K\bigl(g(Y,Z)X-g(X,Z)Y\bigr),
\]
where \(K\) is constant. The case \(K=0\) retains flatness. The source does not specify whether it asks for nonzero \(K\), which other Frobenius axioms should be retained or modified, or what properties the proposed replacement must preserve. Thus this question specifies a research target, not a uniquely defined generalized structure or conjecture.

The supplied motivation is qualified: certain Frobenius manifolds \emph{seem} to yield geometric Poisson brackets of order \(3\), with the WDVV equation singled out by the parenthetical statement that it does so. The record supplies neither the construction nor the meaning of geometric Poisson bracket needed to turn this motivation into an additional formal assertion.
\subsection{Short English statement}
The problem asks whether Frobenius structures admit a projectively invariant formulation and an extension using metrics of constant curvature. The supplied statement leaves the transformation rules and the axioms of the proposed extension unspecified. Its connection with third-order geometric Poisson brackets is supplied as qualified motivation. This is a statement-only draft requiring human mathematical review.
\subsection{Sources}
\begin{itemize}
\item Local UnsolvedMath registry, permanent ID \(801158\), exact source problem code \texttt{catalogue entry}, title ``Nijenhuis Geometry and Integrable Systems --- Problem 11.'' The supplied record attributes the problem to Bolsinov, Konyaev, and Matveev, with proposed year \(2024\), and labels it \texttt{open}; this label is supplied curation, not an independently verified mathematical status.
\item Bolsinov, Konyaev, Matveev, \emph{Research Problems on Nijenhuis Geometry and Integrable Systems} (2024), Problem 11. Supplied source URL: \url{https://arxiv.org/abs/2410.04276}. The background records access on \(2026\text{-}10\text{-}16\) and extraction method \texttt{arxiv\_latex}; these are supplied provenance statements, not claims of access or verification for this draft.
\item Supplied upstream revision: \texttt{d371ee917ac9647fbb06cd6fc529deec6f3dda82}. The registry record was created and updated on \(2026\text{-}10\text{-}16\).
\end{itemize}
\end{document}
